% Part: many-valued-logic
% Chapter: infinite-valued-logics
% Section: goedel

\documentclass[../../../include/open-logic-section]{subfiles}

\begin{document}

\olfileid{mvl}{inf}{god}

\olsection{منطق‌های گودل}

\begin{editorial}
  این بخش پیش‌نویسی بسیار کوتاه دربارهٔ منطق بی‌نهایت‌ارزشی گودل است.
\end{editorial}

\begin{defn}\ollabel{def:goedel} منطق بی‌نهایت‌ارزشی گودل، یعنی
$\LogGod[\infty]$، با ماتریس زیر تعریف می‌شود:
\begin{enumerate}
  \item زبان گزاره‌ای استاندارد~$\Lang L_0$ با~$\lfalse$، $\lnot$،
  $\land$، $\lor$ و~$\lif$.
  \item مجموعهٔ ارزش‌های صدق~$V_\infty$.
  \item $1$ تنها ارزش ممتاز است؛ یعنی~$V^+ = \{1\}$.
  \item توابع صدق با توابع زیر داده می‌شوند:
  \begin{align*}
    \tf{\lfalse} & = 0\\
    \tf{\lnot}[\LogGod](x) & = \begin{cases}
      1 & \text{اگر } x =0\\
      0 & \text{در غیر این صورت}
    \end{cases}\\
    \tf{\land}[\LogGod](x,y) & = \min(x,y)\\
    \tf{\lor}[\LogGod](x,y) & = \max(x,y)\\
    \tf{\lif}[\LogGod](x,y) & = \begin{cases}
      1 & \text{اگر } x \le y\\
      y & \text{در غیر این صورت.}
    \end{cases}
    \end{align*}
\end{enumerate}
منطق~$m$ارزشی گودل نیز به همین شیوه تعریف می‌شود، جز اینکه~$V = V_m$.
\end{defn}

\begin{prop}
  منطق~$\LogGod[3]$ که با~\olref[thr][god]{defn:goedel} تعریف شده،
  همان~$\LogGod[3]$ تعریف‌شده با~\olref{def:goedel} است.
\end{prop}

\begin{proof}
  این امر با مقایسهٔ جدول‌های صدق رابط‌ها در
  \olref[thr][god]{defn:goedel} با جدول‌های صدقی که معادلات
  \olref{def:goedel} تعیین می‌کنند روشن می‌شود:
  \begin{center}
    \begin{tabular}{c|c} 
      $\tf{\lnot}[\LogGod[3]]$ & \\ 
      \hline  
      $1$ & $0$ \\ 
      $1/2$ & $0$ \\
      $0$ & $1$ 
    \end{tabular}
    \quad
    \begin{tabular}{c|ccc} 
      $\tf{\land}[\LogGod]$ & $1$ & $1/2$ & $0$ \\ 
      \hline 
      $1$ & $1$ & $1/2$ & $0$ \\ 
      $1/2$ & $1/2$ & $1/2$ & $0$\\ 
      $0$ & $0$ & $0$ & $0$ 
    \end{tabular}
    \\[2ex]
    \begin{tabular}{c|ccc} 
      $\tf{\lor}[\LogGod]$ & $1$ & $1/2$ & $0$ \\ 
      \hline 
      $1$ & $1$ & $1$ & $1$ \\ 
      $1/2$ & $1$ & $1/2$ & $1/2$ \\
      $0$ & $1$ & $1/2$ & $0$ 
    \end{tabular}
    \quad
    \begin{tabular}{c|ccc} 
      $\tf{\lif}[\LogGod]$ & $1$ & $1/2$ & $0$ \\ 
      \hline 
      $1$ & $1$ & $1/2$ & $0$ \\ 
      $1/2$ & $1$ & $1$ & $0$  \\ 
      $0$ & $1$ & $1$ & $1$ 
    \end{tabular}
  \end{center} 
\end{proof}

\begin{prop}\ollabel{prop:god-infty-m}
  اگر~$\Gamma \Entails[\LogGod[\infty]] !B$، آنگاه برای همهٔ
  $m \ge 2$، $\Gamma \Entails[\LogGod[m]] !B$.
\end{prop}

\begin{proof}
  تمرین.
\end{proof}

\begin{prob}
  \olref[mvl][inf][god]{prop:god-infty-m} را ثابت کنید.
\end{prob}

در واقع، عکس این گزاره نیز برقرار است.

مانند~$\LogGod[3]$، در~$\LogGod[\infty]$ همهٔ~!!{formula}های معتبر از
دید شهودگرایانه همان‌گویی‌اند، و همان مثال‌های ناهـمان‌گویی در
$\LogGod[\infty]$ نیز همان‌گویی نیستند:
\begin{align*}
  & p \lor \lnot p && (p \lif q) \lif (\lnot p \lor q) \\
  & \lnot\lnot p \lif p && \lnot(\lnot p \land \lnot q) \lif (p \lor q) \\
  & ((p \lif q) \lif p) \lif p && \lnot(p \lif q) \lif (p \land \lnot q)
\end{align*}
مثال~$(p \lif q) \lor (q \lif p)$ که~!!{formula}ای نامعتبر از دید
شهودگرایانه اما همان‌گویی~$\LogGod[3]$ است، در~$\LogGod[\infty]$ نیز
همان‌گویی است. در واقع، می‌توان~$\LogGod[\infty]$ را منطق
شهودگرایانه‌ای مشخص کرد که طرح‌وارهٔ
$(!A \lif !B) \lor (!B \lif !A)$ به آن افزوده شده است. این نتیجه را
مایکل دامت نشان داد؛ ازاین‌رو~$\LogGod[\infty]$ را اغلب منطق
گودل--دامت~$\Log{LC}$ می‌نامند.

\begin{prob}
  نشان دهید~$(p \lif q) \lor (q \lif p)$ همان‌گویی
  $\LogGod[\infty]$ است.
\end{prob}

\begin{prob}
  نشان دهید~$(p \lif q) \lor (q \lif r) \lor (r \lif s)$، که
  همان‌گویی~$\LogGod[3]$ است، در~$\LogGod[\infty]$ همان‌گویی نیست.
\end{prob}

\end{document}
